% !PRINTED 6
\weberpage{6}

\gesperrt{5. Groups which have this property are called commutative groups or also Abelian groups.}

If the elements of two groups
\[
a, b, c, d \ldots
\]
and
\[
a', b', c', d' \ldots
\]
correspond to one another in a mutually one-to-one manner, so that whenever $ab = c$, then also $a'b' = c'$, the groups are called \gesperrt{isomorphic}, and the evident theorem holds that two groups isomorphic with a third are isomorphic with one another. Accordingly one may gather all groups that are isomorphic with one another into a class of groups, which is itself again a group, whose elements are the generic concepts obtained when one gathers the corresponding elements of the individual isomorphic groups into one general concept. The individual mutually isomorphic groups are then to be regarded as different representatives of one generic concept.

The properties of groups that can come into consideration are of different kinds. They may, namely, either attach to the particular groups and be derived from the nature of the elements of which the group consists, or else from the nature of the law of composition. Or they may attach to the groups as such, and must then admit of being derived solely from the definition of the concept of a group. The latter properties are common to all isomorphic groups and may be designated as invariant properties of the group. If a comparison is permitted, one might recall the distinction between the metric and the projective properties in geometry.

To the first kind of properties, which are derived from the particular nature of the elements, belong, e.\,g., in the case of permutation groups, the properties of transitivity and intransitivity, of primitivity and imprimitivity; to the invariant properties belong permutability or non-permutability, the degree, the divisors and their index, the normal divisors. It is with these invariant properties, for which it is immaterial from which representative of a class of isomorphic groups they are derived, that we have first to occupy ourselves.
